\documentclass[../main.tex]{subfiles}
\begin{document}
\paragraph{Katyusha is not really a learning rule but an update rule, the main focus of the
rule is how the calculated gradients are used to update the weights and biases
to accelerate Stochastic Gradient Descent. \\
Since my implementation is based around full-batch gradient descent the stochastic part of the katyusha algortithm is changed to work with full batches.}
\subsection{Main Elements of the Rule}
\begin{itemize}
  \item The rule uses the Nesterov’s momentum to boost the descent of gradient but also adds a negative momentum to keep it anchored.
  \item There are two loops in this rule. The outer one is for $epochs$ and the inner one runs for $(2 \times \text{Number of data samples}) = m$. For our case the inner loop runs for 3 iterations
  \item Four parameters are calculated before starting the training.
    \begin{itemize}
            \item $\tau_{1}$ from the nestrov momentum $ = min{\left(\frac{\sqrt{m\sigma}}{\sqrt{3L}}, \frac{1}{2}\right)}$ or $\frac{2}{s+4}$ in our case where $s$ is the epoch number.
\item $\sigma$ is convexity constant. $0$ in our case.
\item $L$ is smoothness constant. Calculated by finding the eigenvalue of $X^{T}X$
\item $\tau_{2}$ is katyusha momentum $ = \frac{1}{2}$ this can be tuned but default is good enough.
\item $\alpha$ is the learning rate $ = \frac{1}{3\tau_{1}L}$
    \end{itemize}
      \item There are three main vectors:
    \begin{itemize}
      \item $\tilde{x}$ this is the inintial snapshot with the model wights and biases. It is updated at the end of each epoch and it the trained model at the end of training.
      \item $y^{k}$ this is used to hold gradient during the inner iterations at each epoch.
      \item $z^{k}$ this is used to hold the past gradient information to build a momentum.
    \end{itemize}
    \begin{itemize}
  \item After every $m$ iterations we take a snapshot gradient of the whole dataset.
  \item In the inner loop we calculate
    \begin{itemize}
      \item $x_{k+1} = \tau_{1}z_{k} + \tau_{2}\tilde{x}^{s} + (1 - \tau_{1} - \tau_{2})*y_{k}$
      \item $\bigtriangledown_{k+1} = \bigtriangledown_{f(x_{k+1})}$ because I am using full batch.
        \item $z_{k+1} = z_{k} - \alpha * \bigtriangledown_{k+1}$
        \item $y_{k+1} = x_{k+1} - (\frac{1}{3L})*\bigtriangledown_{k+1}$
    \end{itemize}
  \item final snapshot $\frac{1}{m}\sum_{j=0}^{m-1}{y_{sm+j+1}}$
\end{itemize}
\end{itemize}
\subsection{Insights}
\begin{itemize}
  \item The rule takes much more time than just gradient descent.
    \item The rule is an optimizer like ADAM.
    \item The rule drops the loss rapidly at first but then slows down as it reaches near $0$.
    \item The rule can work with very large networks which the normal full batch gradient descent failed to do.
\end{itemize}
\end{document}
